% generated by scripts/paper_framing.py, rule framing-v2 (chosen after reading, see DECISIONS 2026-11-23T20:30Z); do not edit
We introduce a procedural benchmark of drawn directed graphs with exact labels, a verifier that accepts every optimal repair edit, and answer-flipping twins that cap any image-blind predictor at $0.50$ per task, and use it to test whether training a small vision-language model only to repair drawings improves its answers to reachability and edge-criticality questions more than training it on those questions directly. On an adapted $500$-graph VGCure evaluation (drawings by other authors in a style none of our adapters saw, with our prompts, images scaled to our input size and strict scoring), the repair-trained adapter (arm~D) answers edge-existence and reachability questions better than direct question-answer fine-tuning (arm~B): the difference is $+0.056$ (95\% interval [$+0.038$, $+0.075$], mean of 3 seeds, above zero at every seed). On our own pre-registered locked test, which decides the hypothesis, D does worse than B: D$-$B is $-0.105$ (95\% interval [$-0.117$, $-0.093$]), so the hypothesis is not supported at this scale. The loss is confined to edge criticality, a task VGCure does not contain, where D answers \emph{true} to 79\% of the questions (true is right for half); the VGCure gain comes mostly from false edge claims and unreachable pairs, which D rejects more often than B. A parse-then-solve arm, which reads the graph with a separately trained extraction adapter, performs best on our benchmark (E$-$B $+0.065$) but scores $0.092$ below direct question-answer fine-tuning on VGCure, where its parser gives no usable graph for 7\% of the images. This is a pilot-scale study: all adapters were trained on $512$ synthetic root graphs with one 2B base model (Qwen3.5-2B, LoRA), and nothing here measures real diagrams.
